\bookchapter{Chapter 32: Phone Safety Emergency Handbook}

\chapterintro{This chapter is an emergency plan. We hope you never need it — but if your phone is lost, stolen, or used in a scam, it will help you cut the damage. In earthquakes, rainstorms, and other disasters, a phone is also how you get warnings, share your location, and stay in touch. Read it through now, screenshot the key steps or print it and keep it at home — because in an emergency your phone may already be gone, and in a disaster it may be dead or offline.}

\booksection{Scenario 0: Earthquake / extreme weather — how your phone helps you stay alive}

This is about personal safety, not accounts and money. For warning setup, see Chapter 13, “Weather alerts, disaster warnings, and earthquake early warnings.” Here we only cover what to do after the alarm sounds.

\textbf{Earthquake alarm: Drop, cover, hold — don't just film it}\par

\begin{enumerate}
  \item Take cover nearby: Drop, cover, hold on (under a table, inside corners of load-bearing walls, next to sturdy furniture). Protect head and neck. Stay away from glass, exterior walls, and hanging lights.
  \item Don't blindly rush downstairs: The main shock often lasts only seconds to tens of seconds; stairwells can be more dangerous. Never use elevators.
  \item Evacuate in order after shaking stops: Use the stairs, watch for aftershocks and falling objects. Reach open ground away from poles, billboards, and unstable buildings.
  \item Use your phone wisely: Short messages to emergency contacts ("I'm at XX, safe/hurt, need help") save more battery than long calls. Prefer text/SMS/WeChat when you can.
  \item No signal? See Chapter 7 on satellite SOS / satellite messaging (on supported models), or try higher/open ground. In some cases, network-free Bluetooth walkie-talkie apps help too.
  \item Turn off non-essential features: In an emergency, toggle airplane mode on a timer, turn off Bluetooth scanning, lower brightness — stretch battery life by hours. Save the last bar for "I'm safe" and rescue messages.
\end{enumerate}

\begin{tcbwarn}{Rainstorm / typhoon / flash flood: Act during the warning stage}
\begin{itemize}
  \item On a red alert: Stay indoors as much as possible. Avoid underground spaces, underpasses, riverbanks, and steep slopes.
  \item Trapped outdoors: Move to higher floors of sturdy buildings. Don't wade through floodwater (missing manhole covers and live wires). When calling 119/110, give landmarks and water level.
  \item Before traveling: Check weather apps and local alert accounts. On long drives, avoid mountain roads along water and cliffs.
  \item Family plan: Agree that "if we lose contact, use the last WeChat/SMS location" so nobody burns battery on endless calls.
\end{itemize}

\end{tcbwarn}

\begin{tcbpractice}{Four things to do before you need them}
\begin{enumerate}
  \item Turn on system earthquake warnings / a trusted warning app, and let it break through silent mode.
  \item Pin 2–3 emergency contacts at the top of Contacts and show them on the lock screen (iOS/Android Medical ID or lock-screen message).
  \item Keep power banks at home and in your bag. Charge up and pack a cable before long trips.
  \item Memorize at least one number you can dial without Contacts (a family mobile). Write it on paper in your wallet.
\end{enumerate}

\end{tcbpractice}

\booksection{Scenario 1: Phone lost or stolen — the golden 30 minutes}

The first 30 minutes after losing a phone are “golden time.” Handle them right and you protect your money and privacy. Follow this order and skip nothing:

\textbf{Standard 30-minute playbook}\par

\begin{enumerate}
  \item Borrow a phone and call your own number (minute 1): It may just be nearby; a kind stranger can return it.
  \item Lock the phone remotely (minutes 2–5): On another phone or computer, sign in to Find My:
  \item Android brands: Huawei cloud.huawei.com / Xiaomi i.mi.com / OPPO cloud.oppo.com, etc.
  \item Suspend the SIM (minutes 5–10): Call your carrier (China Mobile 10086 / Unicom 10010 / Telecom 10000) and report the phone lost to suspend the SIM. This is the most critical step — if scammers get your SIM, they receive your codes and can reset every password.
  \item Freeze payment accounts (minutes 10–20):
  \item WeChat: Report to police at 110, then contact WeChat support, or WeChat Security Center > Freeze account
  \item Banking apps: Call the number on the back of your bank card and ask to freeze online banking
  \item Tell family and friends (minutes 20–25): Post on Moments or have family spread the word that your phone is lost, so scammers can't use your identity to trick others.
  \item Report to police (minutes 25–30): Call 110 or go to the nearest station and get a report receipt. You'll need it for insurance and bank disputes later.
\end{enumerate}

\booksection{Scenario 2: You got scammed — the 72-hour window}

If you realize you've been scammed, staying calm is critical. Panic leads to bad decisions and more losses. Do this:

\begin{enumerate}
  \item Stop all transfers immediately: No matter what they say, don't send another cent. "Send one more deposit and we'll refund everything" is the classic second-round scam.
  \item Keep every piece of evidence: Don't delete chats! Don't delete call logs! Screenshot chats, transfers, and the other party's account details. These are key for reports and getting money back.
  \item Call 96110 right away: China's national anti-fraud hotline. Operators guide emergency stop-payment — if the transfer was recent (usually within 24 hours), the bank may freeze the other account and intercept funds.
  \item Call 110: Bring your phone, ID, and transfer receipts to the nearest station for a statement.
  \item Contact your bank/payment platform: Report the scam and ask to freeze the accounts involved. For credit-card fraud, ask the bank to start a chargeback.
  \item Change every password: On another safe device, change passwords for Apple ID/brand account, WeChat, Alipay, and banking apps, and turn on two-factor authentication.
\end{enumerate}

\booksection{Scenario 3: Phone infected with malware — what now?}

If your phone shows strange pop-ups, installs weird apps on its own, or your bill/data spikes — it's likely malware:

\begin{itemize}
  \item Step 1: Enter safe mode — on Android, power off with the power button, then power on and hold volume-down after the vibration until "Safe mode" appears. Third-party apps are disabled in safe mode, making cleanup easier.
  \item Step 2: Uninstall suspicious apps — in safe mode, Settings > Apps, sort by date, remove unknown recent installs.
  \item Step 3: Scan with a security app — install "Tencent Mobile Manager" or "360 Mobile Security" and run a full scan.
  \item Step 4: Factory reset (last resort) — if nothing works, back up important data and factory reset. Warning: A factory reset erases everything. Back up photos, contacts, and key files first.
  \item iPhones almost never "get viruses" — iOS's sandbox makes malware hard to survive. If your iPhone acts up, it's usually an app bug or a settings issue. A factory reset fixes 99\% of problems.
\end{itemize}

\booksection{Scenario 4: Account hijacked — take back control fast}

\begin{enumerate}
  \item Change your password now: If you can still sign in, change it immediately. Make it completely different from the old one.
  \item Kick other devices: In account security settings, find "Login devices" and sign out anything you don't recognize.
  \item Check login history: Review the last 7 days. Strange locations mean the account really was taken.
  \item Turn on two-factor authentication: If you haven't, do it now. It's the most effective way to prevent another hijack.
  \item Check linked info: Verify the phone number and email on the account weren't changed. If they were, use account recovery to take it back.
  \item Check payment settings: If WeChat/Alipay was hit, look for altered password-free payments or auto-debits, and bank cards added to unfamiliar accounts.
\end{enumerate}

\booksection{Scenario 5: Blackmail message — don't panic, don't pay}

Two common scripts: (1) “I have your private video/photos; pay or I'll send them to everyone in your contacts”; (2) “Your phone is hacked; pay or I'll lock it.” If this happens:

\begin{itemize}
  \item 99.9\% is a scam: They don't have anything on you. They send the same message to tens of thousands of people and wait for someone to pay.
  \item Don't reply, don't pay: Replying tells them "this number is live" and marks you as a target for more harassment.
  \item Screenshot the evidence, then block them.
  \item Change passwords: If it eases your mind, change passwords on your main accounts.
  \item If there's a real privacy leak risk: Report to police. They have experience with these cases.
\end{itemize}

\booksection{Scenario 6: Deepfake video / voice cloning scams — seeing is no longer believing}

Since 2024, deepface and voice cloning let scammers “become” your family, boss, or friends on video calls. “Mom, I'm in trouble, send money” — the person looks and sounds exactly like your son, but it's fake. This is one of the hardest and most dangerous scams right now.

\textbf{Step 1: Spot the tells — four common giveaways in AI fakes}\par

\begin{itemize}
  \item Abnormal blinking: Real people blink 15–20 times a minute; AI faces often blink too little or with stiff rhythm.
  \item Lips out of sync: Watch the mouth against the audio — there's often a fraction-of-a-second lag.
  \item Blurry face edges: Hairlines, ears, and hair edges may show "fuzzy" fringes or flicker.
  \item Inconsistent lighting: Highlights and shadows on the face don't match the background light, or features "slide" wrong when the head turns.
\end{itemize}

\begin{tcbinfo}{Step 2: Three ways to check if they're real}
\begin{itemize}
  \item Ask for a specific action: Turn the head, cover the face, block the mouth then open and close it. Deepfakes break on occlusion and big head turns — faces warp or lag behind motion.
  \item Ask only they would know: Invent a fake question like "which clinic did we take the dog to last time?" Scammers can't answer or will stall.
  \item Call back on a known number: Hang up and dial the number saved in your contacts. Scammers can't control the real phone number — the truth comes out immediately.
\end{itemize}

\end{tcbinfo}

\begin{tcbwarn}{Step 3: If you already lost money — report it}
\begin{enumerate}
  \item Stop transfers now: No matter how "urgent," pause. Anyone pushing "transfer fast, don't hang up" is almost certainly a scammer.
  \item Keep every piece of evidence: Screen-record and screenshot video call frames, chats, their account, transfer receipts — key material for technical investigation.
  \item Call 96110: Say it's a "deepfake voice/video" scam; operators can start emergency stop-payment.
  \item Call 110: Bring all evidence to the station and explain AI forgery was used — this helps police classify the case.
  \item Warn family and friends: Scammers may use your fake video on others. Post a heads-up in Moments or the family group chat.
\end{enumerate}

\end{tcbwarn}

\booksection{Scenario 7: SIM swap — your number gets “stolen”}

Many people haven't heard of this attack, but it's devastating. A SIM swap means an attacker fakes your identity and gets a replacement SIM at a store. Once that succeeds, your phone suddenly loses service, while their phone with the new SIM receives every SMS code — Alipay, banks, WeChat all wide open.

\begin{tcbinfo}{How to notice: Three warning signs of a SIM swap}
\begin{itemize}
  \item Phone suddenly shows "No service" / "No SIM" even though you didn't touch the card and restarting doesn't help — often the moment of the swap.
  \item Carrier texts say "SIM replacement successful" or "service change successful" — but you never requested it.
  \item You suddenly stop getting verification codes — and other texts too, because SMS has been hijacked to their card.
\end{itemize}

\end{tcbinfo}

\begin{tcbwarn}{Protection — put a lock on your SIM too}
\begin{itemize}
  \item Set a SIM PIN: Path Settings > Phone > SIM PIN. After enabling, every restart or SIM swap needs the PIN. Even if they get the card, they can't use it.
  \item Enable carrier two-step verification: In China Mobile/Unicom/Telecom apps, find "Security Center" or "Service verification," turn on replacement-SIM alerts and a second password for service changes. Then staff will demand extra ID checks if someone tries to replace your SIM.
  \item Don't rely on SMS for important accounts: For banks, Apple ID, and other high-value accounts, prefer authenticator apps (Google Authenticator, Microsoft Authenticator) or hardware keys over SMS codes.
  \item Watermark ID copies: When a business needs a copy, write "For XX use only" — attackers often impersonate you with leaked ID info.
\end{itemize}

\end{tcbwarn}

\textbf{If it already happened — speed is everything}\par

\begin{enumerate}
  \item Call your carrier immediately to suspend: Mobile 10086 / Unicom 10010 / Telecom 10000. Freeze SIM replacement and stop the attacker from using your number.
  \item Sign in to key accounts from another device: On a computer or family phone, change passwords for banks, Alipay, WeChat, Apple ID and kick unknown devices.
  \item Freeze payment accounts: Alipay 95188 report loss, WeChat Security Center freeze, bank support freeze online banking.
  \item Go to a carrier store with your ID: Get a new SIM in person and restore your number.
  \item Report to police: SIM swap counts as "infringing citizens' personal information." Bring evidence (suspension time records, financial loss, replacement SIM records).
  \item Full audit: After getting your number back, check recent logins and transactions for misuse.
\end{enumerate}

\begin{tcbwarn}{Prevention beats cure — five daily habits}
\begin{enumerate}
  \item Lock screen passcode + two-factor authentication — both on
  \item "Find My Phone" confirmed on
  \item SIM PIN set (Settings > Phone > SIM PIN) so a lost phone can't use your card
  \item Don't install shady apps; don't tap unknown links
  \item Back up important data regularly (photos, contacts, chat history)
\end{enumerate}

\end{tcbwarn}

\begin{tcbpractice}{Hands-on: Make your own emergency contact card}
\begin{enumerate}
  \item On a piece of paper, write: carrier support (Mobile 10086 / Unicom 10010 / Telecom 10000), bank support, 96110 anti-fraud hotline, 110 police, family numbers.
  \item Write the registration email and backup contact for your Apple ID / brand account.
  \item Keep that paper in your wallet or on the fridge — not only on your phone, because in an emergency your phone may be gone.
  \item Check now: Is "Find My Phone" on? Is the SIM PIN set?
\end{enumerate}

\end{tcbpractice}

\begin{tcbsummary}{Chapter recap}
\begin{itemize}
  \item Lost phone golden 30 minutes: Borrow a phone lock remotely suspend SIM freeze payments warn contacts report to police
  \item Scammed 72-hour window: Stop transfers keep evidence call 96110 report contact bank change passwords
  \item Malware: Safe mode uninstall suspicious apps security scan factory reset (last resort)
  \item Hijacked account: Change password kick devices check logins enable 2FA check linked info check payments
  \item Blackmail messages are 99.9\% scams — don't reply, don't pay, screenshot and block
  \item Deepfake voice/video scams: Watch blinking, lips, edges, lighting; ask them to turn their head and cover their face; ask private questions; call back on a known number
  \item SIM swap: Treat sudden "No service" as a red flag. Set a SIM PIN, enable carrier two-step verification, use authenticator apps for important accounts
  \item Five daily habits: Lock screen + 2FA, Find My Phone, SIM PIN, no shady apps, regular backups
\end{itemize}

\end{tcbsummary}
